\documentclass[11pt]{letter}
\usepackage[a4paper,margin=2.6cm]{geometry}
\usepackage[T1]{fontenc}
\usepackage[utf8]{inputenc}
\signature{D\'aniel Bachrathy \\ (on behalf of the authors) \\[4pt]
Department of Applied Mechanics \\
Faculty of Mechanical Engineering \\
Budapest University of Technology and Economics \\
M\H{u}egyetem rkp.\ 3., H-1111 Budapest, Hungary \\
bachrathy@mm.bme.hu}
\address{Department of Applied Mechanics \\
Budapest University of Technology and Economics \\
Budapest, Hungary}
\date{\today}

\begin{document}
\begin{letter}{Editorial Office \\ Journal of Sound and Vibration}

\opening{Dear Editors,}

Please consider our manuscript, \emph{``Multiplication-Free Stochastic
Semi-Discretization for the Efficient Moment Stability Analysis of Vibration
Systems with Delay''} by D\'aniel Bachrathy and Henrik T.~Sykora, for
publication in the \emph{Journal of Sound and Vibration} as a regular research
paper.

The moment stability analysis of stochastic delay differential equations ---
the natural model class for machining dynamics, delayed feedback control, and
connected vehicles under noise --- has long been limited by the
$\mathcal{O}(p^4)$ cost of assembling and multiplying the second-moment step
operators of the semi-discretization framework. The manuscript removes this
bottleneck and, we believe, changes what is practically computable in this
field:

\begin{enumerate}
\item A \textbf{multiplication-free reformulation} evaluates the identical
      discrete second-moment operator in strictly $\mathcal{O}(p^2)$ time.
      Measured on the same machine and the same operator, the advantage over
      the explicit period product reaches $1.3\times10^{3}$ in wall-clock time
      and $6.3\times10^{3}$ in memory allocation already at $p=192$, beyond
      which the classical route is infeasible while the new one continues to
      $p=4096$.
\item A \textbf{Kronecker-factored It\^o operator} removes the
      $\mathcal{O}(d^4)$ state-dimension wall: second-moment computations for
      systems with one hundred degrees of freedom --- covariance eigenproblems
      with millions of unknowns --- complete within minutes on a workstation,
      with a GPU port included in the accompanying open-source package.
\item \textbf{High-order collocation covariance blocks} raise the measured
      second-moment convergence order to $2S$ (rates consistent with order
      eight at $S=4$), where all previous schemes of this family were limited
      to first order.
\item The methodology is demonstrated on a \textbf{two-degree-of-freedom
      milling process with spindle-speed variation}, time-varying delay, and
      multiplicative cutting-force noise. The resulting process map separates,
      for the first time at this fidelity, the deterministic stability gain of
      spindle-speed variation from the part that stochasticity takes back, and
      shows that the practically usable window is set by a surface-quality
      (stationary-variance) boundary lying strictly below all stability
      boundaries. The stationary variance is validated against an independent
      Monte-Carlo simulation to within $1.4\%$.
\end{enumerate}

All results are reproducible: the methods are released in an open-source
Julia package, every figure has a generating script, and the numerical
experiments specify hardware, solver settings, and per-figure computational
cost.

The manuscript is original, has not been published previously, and is not
under consideration elsewhere. All authors have approved the submission and
declare no competing interests. Given the topic, the readership of the
\emph{Journal of Sound and Vibration} --- at the intersection of nonlinear
vibrations, machining dynamics, and computational methods --- is the natural
audience for this work.

Thank you for your consideration.

\closing{Yours sincerely,}

\end{letter}
\end{document}
